\section*{Abstract}
\phantomsection % TOC link target; \section* creates none
\addcontentsline{toc}{section}{Abstract}

Human spaceflight missions expose astronauts to complex physiological, environmental, and operational stressors that may compromise their health, well-being, and performance. Continuous health monitoring is therefore essential for identifying potentially adverse changes and supporting informed decision-making throughout space missions. However, the integration and interpretation of heterogeneous health-related data remain significant challenges.

This work presents \textbf{Astronaut Health Intelligence (AHI)}, a predictive health monitoring platform that moves from health monitoring to predictive health protection. AHI assumes as its data sources the equipment in use on the International Space Station: wearables such as the Bio-Monitor and the Actiwatch Spectrum, medical equipment such as the blood pressure and ECG monitor and the Ultrasound 2, the optical coherence tomography and the hearing assessment used for periodic eye and ear examinations, the CEVIS, ARED and T2 exercise devices, and the environmental monitors for air quality, carbon dioxide, noise and radiation. It organizes health into ten dimensions, and uses the five hazards of human spaceflight defined by NASA as the axis of its risk analysis. It follows the career radiation dose of each astronaut against the limit published in NASA-STD-3001, and it recommends preparing the radiation shelter when space weather notifications or the on board monitors indicate a solar particle event.

The analysis works on three levels. Real-time monitoring compares each astronaut with their own baseline and detects acute changes, slow drifts and missing data. Predictive analytics projects trends and estimates risks, starting with transparent statistical rules and adding machine learning models such as Isolation Forest and XGBoost as real data allow. Prescriptive intelligence recommends exercise, nutrition, sleep, well-being and medical follow-up actions, each with its evidence and an indication of whether the crew can act on board. A Large Language Model (LLM) writes the statistical predictions in plain language under strict guardrails and never computes them, and a health assistant answers questions about the data by text and voice. Internet access is useful but not essential: the demonstrator is a web application that uses the Claude API to simulate the AI language features, and, because a real system on board cannot depend on the internet, the architecture includes an on board backend and a contingency plan that keep monitoring, alerts and recommendations working without a connection, with a local model specialized in AHI for the language features. The platform is designed for a global and diverse crew, with the interface in English, Portuguese, Spanish, Japanese and German, role-based access for crew members and flight surgeons, and audio description for blind and low vision users, for people with dyslexia and for astronauts whose eyes and hands are busy with a task. Public NASA data from the Open Science Data Repository and related sources ground and test the analysis.

AHI is designed as a decision-support tool. It does not diagnose, does not prescribe, and does not replace clinical evaluation or qualified medical personnel.

\noindent\textbf{Keywords:} Astronaut Health Monitoring; Predictive Analytics; Anomaly Detection; Space Medicine; Artificial Intelligence; Human Spaceflight; International Space Station.
